% !TeX root = ../main.tex
\section{Review baseline and integration priorities}
\label{sec:tracker-review}
The selected sensor placements, module populations, stereo frames, support
datums and fixed service boundaries form the baseline for review. Its intent is
to provide a common, reproducible starting point for geometry, engineering and
tracking experts. A successful software build or a merged design contribution
does not close a failed engineering screen or establish human design sign-off.

The local models have substantial construction and transport evidence. The
remaining work is concentrated at coupled interfaces, where independently
reasonable subsystems can still conflict. The review should address
\cref{tab:tracker-review} before freezing an integrated detector release.

\begin{table}[!htbp]
\centering\small
\caption{Cross-system review priorities and the evidence needed to close them.
Existing failed cases remain controls for subsequent revisions.}
\label{tab:tracker-review}
\begin{tabular}{@{}>{\raggedright\arraybackslash}p{34mm}>{\raggedright\arraybackslash}p{106mm}@{}}\toprule
Interface & Required review outcome\\\midrule
Acceptance and transitions & Combined finite-plane hit maps and hit multiplicity
over the luminous region, both charges and the specified momentum range; retain
old-datum and original-annulus losses. Establish efficiency with material,
response and reconstruction.\\
Common mechanics & Qualify global cages, ring joints, petal keys and thermal
motion. Confirm the 17-station long-stave recommendation with measured elastic
properties, joint compliance and global-frame loads. Provide handling and
installation tooling, especially for long-strip petals.\\
Cooling and thermal contacts & Couple realistic electronics/leakage power,
anisotropic spreading and contact resistance to pressure drop, flow sharing and
two-phase stability. Resolve short-petal thermal failures and the long-petal
lateral-path screen, retaining warm and adverse cases.\\
Power, data and routing & Specify actual sensors/ASICs, harnesses, insulation,
terminations and bandwidth. Resolve short-strip nominal trunk overfill and
radial turns, long-strip adverse packing and warm return drop, and pixel
trunk/flange packing. Provide strain relief and access within fixed limits.\\
Integrated material model & Replace duplicated standalone service cells once;
audit all shared volumes and interfaces together. Produce full directional
radiation/interaction-length maps and a constituent inventory before quoting a
tracker material budget or mass.\\
Measurement and alignment & Preserve true sensor axes and IDs, qualify strip
pair transport to a common plane, and define response, survey and alignment
covariances. Validate conversion/navigation before interpreting fitted-track
resolution.\\\bottomrule
\end{tabular}
\end{table}

The immediate recommendation is to review the selected layouts and local load
paths together, then resolve common services and barrel/endcap transitions in
one integrated model. This order prevents a subsystem fix from consuming space
already allocated to another subsystem. Fixed envelopes and surviving sensor
identities should remain comparison controls; proposed revisions must report
their material, acceptance and engineering consequences explicitly.

The report sources retain hashed data snapshots, exact input/producer
provenance, actual native execution identities and reproducible figure recipes.
Compilation uses checked-in vector figures and tables only. Regeneration does
not rerun simulation or change geometry. The reviewed publication snapshot is
therefore distinct from the native validation state and from any future
accepted detector release.
